
FastText content classification forms the basis of the DCLM-BASELINE filtering recipe \citep{Li2024DCLM} and is widely used in pretraining data curation. The classifier assigns quality scores based on similarity to Wikipedia and OpenWebText reference corpora. Because these reference corpora contain systematic patterns of demographic-occupation co-occurrence (e.g., formal descriptions of occupations alongside demographic markers), filtering based on similarity to these corpora may alter not only document quality but also the demographic association structure of the retained text.

Prior work on pretraining data curation has focused primarily on downstream task performance. DCLM \citep{Li2024DCLM} demonstrates that fastText filtering achieves 64\% MMLU at 7B scale with 2.6T tokens, but explicitly excludes bias and safety from its evaluation scope. FineWeb \citep{Penedo2024FineWeb} and DoReMi \citep{Xie2023DoReMi} similarly evaluate curation decisions by performance metrics alone. Dolma \citep{Soldaini2024Dolma} documents demographic limitations of web corpora but does not quantify how filtering operations transform demographic structure. On the fairness side, benchmarks such as BBQ \citep{Parrish2022BBQ}, WinoBias \citep{Zhao2018WinoBias}, and StereoSet \citep{Nadeem2021StereoSet} measure model output disparities but do not trace them to specific pretraining data curation decisions.

This paper addresses the gap between these two lines of work by measuring how fastText quality filtering alters corpus-level demographic-occupation association structure. The investigation proceeds through a series of sub-hypotheses within the Path-Dependent Curation Fairness Hypothesis (PCFH) framework:

\begin{itemize}
\item \textbf{H-E1 (Existence):} Whether fastText filtering creates systematic, monotonic changes in conditional entropy H(occupation|demographic). Gate: at least 5\% relative entropy change between extreme configurations.
\item \textbf{H-M1 (Mechanism, corpus level):} Whether conditional log-odds of demographic-occupation co-occurrences correlate with filtering intensity. Gate: Spearman |rho| > 0, p < 0.05.
\item \textbf{H-M2 (Mechanism, model level):} Whether models trained on differently filtered corpora internalize these differences in their logit margins. Gate: Spearman rho > 0, p < 0.01.
\item \textbf{H-M3 (Downstream fairness):} Whether matched-capability models from different curation paths produce distinguishable fairness benchmark outcomes. This sub-hypothesis was not executed.
\end{itemize}

The contributions of this work are: (1) an empirical characterization of how fastText quality filtering alters demographic-occupation conditional entropy and log-odds at corpus level; (2) a model-free corpus audit methodology for measuring these effects prior to model training; and (3) a preliminary, inconclusive pilot of corpus-to-model signal propagation.

